\subsubsection{A formal retraction and purely inseparable cone covers}

\begin{lemma}[Proved: linearization after Frobenius along a formal retraction]
\label{D:formal-retraction-cone-cover}
Let $C$ be a smooth proper connected curve of genus $g$ in a
smooth surface $Y$ over an algebraically closed field $k$ of
characteristic $p>0$.  Write $N=N_{C/Y}$ and $\deg N=-d<0$.
Suppose the entire formal completion $\mathfrak Y=\widehat Y_C$
admits a $k$-morphism
\[
                  r:\mathfrak Y\longrightarrow C,
                  \qquad r|_C=\operatorname{id}_C.
\]
Choose a power $q=p^a$ with $qd>2g-2$, and let
$h:C'=C^{(-a)}\to C$ be the $k$-linear relative Frobenius.
Then
\[
 \mathfrak Y\mathbin{\widehat\times}_C C'
       \simeq\widehat{\operatorname{Tot}(h^*N)}_{C'}
\]
as formal schemes over $C'$, inducing the identity on $C'$
and on its conormal bundle.  The resulting morphism from
this formal line bundle to $\mathfrak Y$ is finite locally
free of degree $q$, is radicial, and has normal ramification
index one.
\end{lemma}

\begin{proof}
We first describe the relative formal-coordinate argument.
Put $I=\mathcal O_Y(-C)$.  Over an affine open of $C$ on
which $N$ is trivial, the retraction splits the map to
$\mathcal O_C$.  A lift $t$ of a generator of $I/I^2$
then identifies the restricted formal algebra with
$\mathcal O_C[[t]]$.  Indeed the induced maps on successive
graded pieces are the isomorphisms
$(I/I^2)^{\otimes n}\simeq I^n/I^{n+1}$; completeness and
separatedness give the claimed isomorphism.  The lifts may
be chosen on affine opens because their nilpotent thickenings
are affine, and the transition maps on global sections are
surjective.  All coordinates here respect the fixed map to
$C$.

Consequently the overlap maps have the form
\[
          t_i=a_{ij}t_j+\sum_{n\geq2}b_{ij,n}t_j^n,
\]
where the coefficients are functions on the curve and
$a_{ij}$ is the transition function of the conormal line.
If the maps have been made linear modulo the $n$th power,
their degree-$n$ errors form a cocycle with values in
\[
 \mathcal Hom(N^{-1},N^{-n})=N^{1-n},\qquad n\geq2.
\]
Replacing a coordinate by itself plus a term of degree $n$
changes that cocycle by a coboundary.  Thus the vanishing
of $H^1(C,N^{1-n})$ permits the next linearization step.
There is no tangent-direction error, because the overlap
maps are maps over the fixed curve.  Nor is there a
first-order obstruction: the retraction already identifies
the first square-zero algebra with
$\mathcal O_C\oplus N^{-1}$.

Now form the indicated base change along $h$.  It is finite
locally free and radicial of degree $q$, since $h$ is.
The pulled-back closed curve is $C'$, its normal bundle is
$N'=h^*N$, and the preceding relative coordinates show that
the base-changed formal scheme is regular.  Its genus is
still $g$ and $\deg N'=-qd$.  For $n\geq2$ we have
\[
               \deg(N'^{\,1-n})=qd(n-1)>2g-2.
\]
Serre duality therefore gives all the required vanishings;
see \cite[Lemmas 53.4.1 and 53.5.2, tags 0BS2 and 0BS6]
{StacksBNegativeCurve}.  The compatible coordinate changes
converge, proving the formal isomorphism.  The ideal of
$C$ pulls back to the ideal of $C'$, so the conormal map
is an isomorphism and the normal ramification index is one.
\end{proof}

\begin{proposition}[Proved: tameness from a formal retraction]
\label{D:formal-retraction-tame}
In the preceding lemma, suppose $Y\to S$ is a projective
resolution of a normal excellent surface germ, isomorphic
off its closed point, whose sole reduced exceptional curve
is $C$.  Then every finite normal generically separable
cover of $S$ that is \'etale off the closed point is tame
at the exceptional valuation.  The same assertion holds
over the strict henselization.  There is no lower bound
on $d$ other than $d>0$.

The formal retraction hypothesis holds, in particular,
if a neighborhood of $C$ admits a morphism $f:Y\to B$
to a smooth curve such that $f|_C:C\to B$ is \'etale.
\end{proposition}

\begin{proof}
Pass to the complete local ring of $S$.  The finite formal
cone map in Lemma~\ref{D:formal-retraction-cone-cover}
algebraizes by \cite[Lemma 30.28.2, tag 09ZT]
{StacksBFiniteCone}.  The proof of
Lemma~\ref{D:finite-formal-cone-cover}, beginning with its
algebraization paragraph, applies verbatim to this already
constructed finite formal algebra: its Stein ring is the
completed section cone of $(C',h^*N)$, is finite over the
completed local ring of $S$, and the actual exceptional
valuation on that cone is the radial one.  This use of the
algebraization argument imposes none of that lemma's
tangent-cohomology hypotheses.

We check the stronger ramification conclusion using the
actual complete DVRs.  At the generic point of $C$, the
retraction supplies the coefficient field $k(C)$ and a
normal parameter $t$.  Before and after the formal base
change their fraction fields are, compatibly,
\[
                 K=k(C)((t))\ \subset\ L=k(C')((t)).
\]
The map on coefficient fields is relative Frobenius of
degree $q$, and $t$ is unchanged.  Thus $L/K$ is purely
inseparable.  These are also the genuine completed DVRs
of the algebraized resolutions, by their identical formal
neighborhoods.

Any given punctured cover of $S$ pulls back to a finite
\'etale cover of the punctured section cone, since the
finite local Stein map has only its vertex above the
closed point.  Proposition~\ref{D:graded-cone-tame} makes
that pullback tame at the radial valuation.  A finite
purely inseparable extension identifies the absolute
Galois groups with their compatible valuation, inertia,
and wild-inertia subgroups, as verified in
Proposition~\ref{D:finite-cone-wild-bound}.  The original
wild inertia therefore acts trivially as well.  This
proves tameness, including separability of the original
residue extension, rather than only a bound on its
ramification index.

For the uncompleted germ and its strict henselization,
the same boundary-DVR comparison as in that proposition
applies, using \cite[Lemma 54.11.1, tag 0BG4]
{StacksBFiniteCone}.  In particular no equivalence of
the two different kinds of completion is assumed.

Finally suppose $f|_C:C\to B$ is \'etale.  It is formally
\'etale by \cite[Lemma 37.8.7, tag 04FF]
{StacksFormalRetractionRoot}.  On every infinitesimal
neighborhood $C_n\subset Y$, the map $C_n\to B$ has a
unique lift $C_n\to C$ extending the identity on $C$.
Formal \'etaleness gives these lifts one square-zero
thickening at a time; uniqueness makes them compatible.
Their inverse limit is the required formal retraction.
\end{proof}

\begin{remark}
A splitting of the tangent-normal sequence alone is
weaker than a retraction of the entire formal
neighborhood.  The latter is the exact extra hypothesis
used here.  The argument uses the Frobenius base change
along that existing retraction and therefore does not
construct one for an arbitrary negative exceptional curve.
\end{remark}
